\section*{Chapter \ref{ch:g4:breathing} --- Breathing}

\begin{solution}{exo:g4:breathing:1}
\omterm{def:g4:breathing:movements}{Breathing} in: the chest expands and air flows in. \omterm{def:g4:breathing:movements}{Breathing}
out: the chest relaxes and air flows out.
\end{solution}

\begin{solution}{exo:g4:breathing:2}
Nose (or mouth), \omterm{def:g4:breathing:organs}{windpipe}, the two \omterm{def:g1:plants-around-us:tree}{branches}, the \omterm{def:g4:breathing:organs}{lungs}, and
finally the millions of tiny air pockets wrapped in \omterm{prop:g4:journey-of-food:absorb}{blood}
vessels.
\end{solution}

\begin{solution}{exo:g4:breathing:3}
About $15$ times a minute. In a sprint it rises to $40$ or
more, deeper too --- the hard-working \omterm{def:g3:bones-and-muscles:muscle}{muscles} need oxygen
faster, so the \omterm{def:g4:breathing:organs}{lungs} load it faster.
\end{solution}

\begin{solution}{exo:g4:breathing:4}
It takes oxygen and returns carbon dioxide. The swap happens
in the \omterm{def:g4:breathing:organs}{lungs}' tiny air pockets, across their thin walls, with
the \omterm{prop:g4:journey-of-food:absorb}{blood}.
\end{solution}

\begin{solution}{exo:g4:breathing:5}
Exhaled air is warmer, moister (the misted mirror), poorer in
oxygen and richer in carbon dioxide.
\end{solution}

\begin{solution}{exo:g4:breathing:6}
Because the \omterm{prop:g4:journey-of-food:absorb}{air-blood} swap needs surface: millions of pocket
walls add up to half a tennis court of meeting surface folded
into the chest. The \omterm{def:g4:journey-of-food:tract}{small intestine} plays the same trick with
its folds for absorbing \omterm{def:g4:journey-of-food:digestion}{nutrients}.
\end{solution}

\begin{solution}{exo:g4:breathing:7}
The nose warms, moistens and filters the air on its way in ---
cold, dry, dusty air reaches the delicate \omterm{def:g4:breathing:organs}{lungs} already
prepared.
\end{solution}

\begin{solution}{exo:g4:breathing:8}
Open answer. Typically about $15$ at rest and $25$--$40$ after
the jumping jacks: the working \omterm{def:g3:bones-and-muscles:muscle}{muscles} demanded more oxygen,
so \omterm{def:g4:breathing:movements}{breathing} sped up and deepened to supply it.
\end{solution}

\begin{solution}{exo:g4:breathing:9}
That the body stores food generously (hours' or days' worth)
but keeps almost no store of oxygen at all --- barely a
minute's. Oxygen must therefore be delivered continuously,
breath by breath, which is why \omterm{def:g4:breathing:movements}{breathing} can never pause long.
\end{solution}

\begin{solution}{exo:g4:breathing:10}
Twenty-five breathers have been taking oxygen from the room's
air and returning carbon dioxide all morning: its air has
grown poorer in one and richer in the other. The prescription
is rule 1: open the windows and air the room.
\end{solution}

\begin{solution}{exo:g4:breathing:11}
Both are the meeting place of \omterm{prop:g4:journey-of-food:absorb}{blood} and oxygen: huge, thin,
blood-lined surfaces where oxygen crosses in and carbon
dioxide crosses out --- gills taking oxygen dissolved in
water, \omterm{def:g4:breathing:organs}{lungs} taking it from air. In air the trout's feathery
gill surfaces collapse and stick together: the great crossing
surface shrinks to almost nothing, and no crossing surface
means no oxygen, however rich the air around it.
\end{solution}
